% ---------------------------------------------------------------------------------------------------------------
% The Atmosphere: Practical Companion -- shared style for every part and every edition.
% The Practical 3 folder keeps an identical copy as preamble.tex (for its slide-deck figure scripts);
% run sync_style.py after editing this file.
% ---------------------------------------------------------------------------------------------------------------
\usepackage[T1]{fontenc}
\usepackage{lmodern}
\usepackage[english]{babel}
\usepackage{microtype}
\usepackage[a4paper,top=2.4cm,bottom=2.4cm,inner=2.5cm,outer=2.2cm]{geometry}
\usepackage{amsmath,amssymb,mathtools}
\usepackage{booktabs,array,tabularx,longtable,colortbl,multirow}
\usepackage{enumitem}
\usepackage{xcolor}
\usepackage{xspace}
\usepackage{needspace}
\usepackage{tikz,pgfplots}
\usepgfplotslibrary{fillbetween,groupplots}
\pgfplotsset{compat=1.17}
\usetikzlibrary{arrows.meta,calc,decorations.markings,decorations.pathreplacing,decorations.pathmorphing,patterns,positioning,shapes.geometric,shapes.symbols,backgrounds,fit,intersections}
\usepackage[most]{tcolorbox}
\usepackage{caption,subcaption}
\usepackage{fancyhdr}
\usepackage{pdflscape}
\usepackage{titlesec}
\usepackage{setspace}
\usepackage{refcount}
\usepackage[round,authoryear]{natbib}
\usepackage{xr-hyper}
\usepackage[hidelinks,pdfusetitle]{hyperref}
\usepackage{url}

\captionsetup{font=small,labelfont=bf,labelsep=period,justification=raggedright,singlelinecheck=false}
\setlist{itemsep=2pt,topsep=4pt,parsep=0pt}
\setlength{\parindent}{0pt}
\setlength{\parskip}{6pt plus 1pt minus 1pt}
\renewcommand{\arraystretch}{1.15}
\numberwithin{equation}{chapter}
\setlength{\bibsep}{4pt}

% ------------------------------------------------------------------ colours (colour-blind friendly) --------------
\definecolor{fblue}{HTML}{0072B2}
\definecolor{fred}{HTML}{D55E00}
\definecolor{fgreen}{HTML}{009E73}
\definecolor{forange}{HTML}{E69F00}
\definecolor{fpurple}{HTML}{CC79A7}
\definecolor{fsky}{HTML}{56B4E9}
\definecolor{fgrey}{HTML}{6B6B6B}
\definecolor{fnavy}{HTML}{1B3A6B}
\definecolor{tempcol}{HTML}{D55E00}
\definecolor{tempcold}{HTML}{0072B2}
\definecolor{fig9label}{HTML}{1B3A6B}
\colorlet{lightblue}{fblue!10}
\colorlet{lightgreen}{fgreen!10}
\colorlet{lightorange}{forange!14}
\colorlet{lightred}{fred!9}

% ------------------------------------------------------------------ units and short-hands -----------------------
\newcommand{\degC}{\ensuremath{^{\circ}\mathrm{C}}\xspace}
\newcommand{\degF}{\ensuremath{^{\circ}\mathrm{F}}\xspace}
\renewcommand{\deg}{\ensuremath{^{\circ}}\xspace}
\newcommand{\gkg}{\ensuremath{\mathrm{g\,kg^{-1}}}\xspace}
\newcommand{\hPa}{\ensuremath{\mathrm{hPa}}\xspace}
\newcommand{\mb}{\ensuremath{\mathrm{mb}}\xspace}
\newcommand{\ms}{\ensuremath{\mathrm{m\,s^{-1}}}\xspace}
\newcommand{\kt}{\ensuremath{\mathrm{kt}}\xspace}
\newcommand{\ws}{\ensuremath{w_{s}}}
\newcommand{\es}{\ensuremath{e_{s}}}
\newcommand{\Vg}{\ensuremath{V_{g}}}
\newcommand{\pgf}{pressure gradient force}
\newcommand{\Td}{\ensuremath{T_{d}}}
\newcommand{\Wm}{\ensuremath{\mathrm{W\,m^{-2}}}\xspace}
\newcommand{\Wmum}{\ensuremath{\mathrm{W\,m^{-2}\,\mu m^{-1}}}\xspace}
\newcommand{\um}{\ensuremath{\mu\mathrm{m}}\xspace}
\newcommand{\umK}{\ensuremath{\mu\mathrm{m\,K}}\xspace}
\newcommand{\Kkm}{\ensuremath{\mathrm{K\,km^{-1}}}\xspace}
\newcommand{\Ckm}{\ensuremath{^{\circ}\mathrm{C\,km^{-1}}}\xspace}
\newcommand{\sigSB}{\ensuremath{\sigma}}
\newcommand{\lmax}{\ensuremath{\lambda_{\max}}}
\newcommand{\Qref}[1]{Question~#1}

% ------------------------------------------------------------------ references across parts -----------------------
% \xref{label} and \xeqref{label}: a reference that may point into another part of the companion. In the full
% companion every label is local. In a single-practical edition, labels of other parts are read from the full
% companion (prefix ext-) and their links open the full companion at that place.
\providecommand{\xref}[1]{\ifcsname r@#1\endcsname\ref{#1}\else\ref{ext-#1}\fi}
\providecommand{\xeqref}[1]{\ifcsname r@#1\endcsname\eqref{#1}\else(\ref{ext-#1})\fi}

% ------------------------------------------------------------------ boxes ---------------------------------------
\tcbset{boxstyle/.style={enhanced,breakable,boxrule=0.6pt,arc=2pt,left=7pt,right=7pt,top=5pt,bottom=5pt,
  fonttitle=\bfseries\sffamily\small,before skip=8pt,after skip=10pt,before=\par\Needspace{6\baselineskip}}}
\newtcolorbox{keybox}[1][Key idea]{boxstyle,colback=lightblue,colframe=fblue,title={#1}}
\newtcolorbox{workedbox}[1][Worked example]{boxstyle,colback=white,colframe=fnavy,coltitle=white,colbacktitle=fnavy,title={#1}}
\newtcolorbox{answerbox}[1][Answer]{boxstyle,colback=lightgreen,colframe=fgreen!80!black,title={#1}}
\newtcolorbox{askbox}[1][A question that comes up]{boxstyle,colback=lightorange,colframe=forange!90!black,title={#1}}
\newtcolorbox{warnbox}[1][Take care]{boxstyle,colback=lightred,colframe=fred,title={#1}}
\newtcolorbox{tryit}[1][Try it yourself]{boxstyle,colback=gray!8,colframe=gray!70,title={#1}}
\newtcolorbox{derivbox}[1][Derivation]{boxstyle,colback=white,colframe=fblue!70!black,borderline={0.4pt}{2pt}{fblue!40,dashed},title={#1}}
\newtcolorbox{questionbox}[1][The question]{boxstyle,colback=fnavy!5,colframe=fnavy!60,title={#1}}

% "where" lists: what each symbol in the equation above stands for, its unit and the value used
\newenvironment{where}{\par\Needspace*{4\baselineskip}\vspace{-2pt}\small\noindent\textit{where}\par\nopagebreak\vspace{1pt}%
  \noindent\tabularx{\linewidth}{@{\hspace{1.2em}}>{$}l<{$}@{\hspace{0.8em}}X@{\hspace{0.8em}}>{\raggedright\arraybackslash}p{0.33\linewidth}@{}}}%
  {\endtabularx\par\vspace{2pt}}

% ------------------------------------------------------------------ headings -----------------------------------
\titleformat{\part}[display]{\centering\bfseries\sffamily\Huge\color{fnavy}}{\partname\ \thepart}{14pt}{\Huge}
\titleformat{\chapter}[display]{\raggedright\bfseries\sffamily\Huge\color{fnavy}}{\chaptertitlename\ \thechapter}{10pt}{\Huge}[\vspace{-2pt}\color{fblue}\titlerule]
\titleformat{\section}{\bfseries\sffamily\Large\color{fnavy}}{\thesection}{0.7em}{}
\titleformat{\subsection}{\bfseries\sffamily\large\color{fnavy}}{\thesubsection}{0.7em}{}
\titleformat{\subsubsection}{\bfseries\sffamily\normalsize\color{fnavy}}{}{0pt}{}
\titlespacing*{\chapter}{0pt}{-20pt}{22pt}

\pagestyle{fancy}
\fancyhf{}
\fancyhead[LE]{\sffamily\small\leftmark}
\fancyhead[RO]{\sffamily\small\rightmark}
\fancyfoot[C]{\sffamily\small\thepage}
\renewcommand{\headrulewidth}{0.3pt}
\fancypagestyle{plain}{\fancyhf{}\fancyfoot[C]{\sffamily\small\thepage}\renewcommand{\headrulewidth}{0pt}}

% ------------------------------------------------------------------ plotting helpers -----------------------------
% Planck curve scaled to its own peak, as a function of wavelength x (micrometres): T is the temperature in kelvin
\pgfmathdeclarefunction{planckn}{2}{%
  \pgfmathparse{(2897.77/(#2*#1))^5*142.32/(exp(4.9651*2897.77/(#2*#1))-1)}}
% Spectral emission on the scale of the Practical 2 data (about 316 times the spectral radiance in W m-2 sr-1 um-1)
\pgfmathdeclarefunction{sheetplanck}{2}{%
  \pgfmathparse{316*1.191042e8/(#1^5*(exp(14387.77/(#1*#2))-1))}}
\pgfplotsset{guide/.style={grid=both,grid style={gray!22},every axis label/.style={font=\small},tick label style={font=\small},
  legend style={font=\footnotesize,draw=gray!50,fill=white,fill opacity=0.92,text opacity=1}}}

% A noon-sun sketch at one site (Practical 2): #1 sun angle (degrees), #2 title, #3 last line.
% The beam is drawn 0.5 cm wide in every sketch, so the sketches can be compared directly.
\newcommand{\sunsketch}[3]{%
\begin{tikzpicture}[line cap=round]
  \pgfmathsetmacro{\sbeta}{#1}
  \pgfmathsetmacro{\strue}{1/sin(\sbeta)}
  \pgfmathsetmacro{\sgl}{0.5*\strue}
  \fill[gray!18] (-0.6,-0.3) rectangle (5.2,0);
  \draw[very thick] (-0.6,0) -- (5.2,0);
  \draw[dashed,gray] (0,0) -- (0,2.9) node[above,font=\scriptsize,text=gray!70!black] {vertical};
  \fill[forange!30,opacity=0.7] ($(0,0)+(\sbeta:2.8)$) -- (0,0) -- (\sgl,0) -- ($(\sgl,0)+(\sbeta:2.8)$) -- cycle;
  \foreach \xx in {0,\sgl}{\draw[-{Stealth[length=2.2mm]},thick,forange!90!black] ($(\xx,0)+(\sbeta:2.8)$) -- (\xx,0);}
  \draw[fred,line width=2.4pt] (0,0) -- (\sgl,0);
  \draw[fblue,very thick] (0.42,0) arc (0:\sbeta:0.42);
  \draw[fgreen!60!black,thick] (90:1.25) arc (90:\sbeta:1.25);
  \pgfmathsetmacro{\zz}{90-\sbeta}
  \node[font=\footnotesize,anchor=north west,align=left,text width=5.4cm] at (-0.6,-0.4)
    {\textbf{#2}\\[1pt]\textcolor{fgreen!45!black}{zenith angle $\theta_z=\pgfmathprintnumber[fixed,precision=1]{\zz}^\circ$} (green arc)\\
     \textcolor{fblue}{sun angle $\beta=#1^\circ$} (blue arc)\\
     \textcolor{fred}{spreading $=\pgfmathprintnumber[fixed,precision=2]{\strue}$ units} (red strip)\\ #3};
\end{tikzpicture}}

% ------------------------------------------------------------------ the weather map frame (Practical 3) -----------
\tikzset{
  mapstate/.style={draw=gray!65,line width=0.25pt},
  isobar/.style={draw=forange!85!black,line width=0.9pt},
  isolabel/.style={font=\fontsize{6.5}{6.5}\selectfont\bfseries,fill=white,inner sep=0.5pt,text=orange!60!black},
  lowmark/.style={font=\fontsize{24}{24}\selectfont\bfseries,text=fred},
  lowval/.style={font=\fontsize{6.5}{6.5}\selectfont\bfseries,text=fred,anchor=west},
  cloudfill/.style={fill=gray!45,opacity=0.35},
  cloudline/.style={draw=gray!65!black,dash pattern=on 3pt off 2pt,line width=0.9pt},
  precipfill/.style={fill=fsky,opacity=0.45},
  streamline/.style={draw=fgreen!85!black,line width=1.1pt},
  exampleline/.style={draw=gray!70,line width=1.5pt},
  fcold/.style={draw=fblue,fill=fblue,line width=1pt},
  fwarm/.style={draw=fred,fill=fred,line width=1pt},
}
% \mapbase{<tikz options>}{<contents drawn on top of the states>}   (the Figure 8 / 10 frame; width 16.4 cm)
\newcommand{\mapbase}[3][]{%
  \begin{tikzpicture}[x=#2,y=-#2,line cap=round,line join=round,#1]
    \clip (170,520) rectangle (990,1150);
    \draw[mapstate] (817.0,1110.0) -- (813.0,1114.7) -- (800.5,1117.0) -- (807.4,1114.1) -- (804.9,1109.9) -- (802.6,1108.5) -- (802.3,1104.2) -- (799.0,1099.4) -- (795.9,1112.2) -- (787.2,1112.7) -- (780.0,1055.2) -- (783.8,939.3) -- (779.9,936.2) -- (859.5,929.7) -- (881.1,1005.2) -- (888.8,1020.8) -- (890.4,1026.9) -- (892.8,1029.5) -- (889.8,1035.1) -- (888.1,1045.5) -- (891.6,1056.9) -- (891.6,1067.5) -- (893.1,1074.1) -- (895.6,1078.3) -- (811.6,1086.7) -- (811.8,1092.0) -- (818.8,1099.8) -- (818.7,1106.1) -- (817.0,1110.0) -- cycle;
\draw[mapstate] (798.3,1116.4) -- (790.3,1117.5) -- (797.0,1115.7) -- (798.3,1116.4) -- cycle;
\draw[mapstate] (731.3,902.6) -- (731.9,905.6) -- (729.7,906.9) -- (729.4,910.1) -- (726.0,912.4) -- (724.7,915.1) -- (725.4,918.4) -- (721.5,925.4) -- (720.5,929.2) -- (721.2,933.3) -- (719.8,936.8) -- (715.4,941.2) -- (715.9,943.4) -- (714.8,945.2) -- (708.3,952.0) -- (707.4,960.6) -- (706.0,964.1) -- (702.1,966.4) -- (700.6,970.5) -- (697.8,973.0) -- (697.0,975.3) -- (697.9,978.3) -- (693.2,981.4) -- (694.0,986.9) -- (689.4,992.8) -- (690.6,995.0) -- (689.3,996.3) -- (691.9,1001.7) -- (691.4,1007.2) -- (693.3,1010.0) -- (692.3,1017.7) -- (600.4,1020.6) -- (600.0,1000.1) -- (595.4,998.8) -- (590.2,999.9) -- (586.2,996.6) -- (587.0,930.9) -- (581.1,888.9) -- (648.7,887.5) -- (716.3,884.6) -- (719.2,889.3) -- (719.3,892.1) -- (710.7,904.1) -- (731.3,902.6) -- cycle;
\draw[mapstate] (671.3,740.3) -- (662.6,731.4) -- (547.2,734.2) -- (544.5,728.8) -- (545.3,723.4) -- (544.9,716.2) -- (543.1,709.9) -- (543.9,708.3) -- (542.1,705.3) -- (542.3,701.8) -- (539.1,698.5) -- (538.4,688.1) -- (532.6,676.3) -- (530.3,665.5) -- (523.8,653.7) -- (528.5,638.6) -- (527.7,635.7) -- (525.7,634.3) -- (526.5,630.7) -- (524.9,624.6) -- (600.2,624.4) -- (671.7,622.3) -- (672.6,626.3) -- (675.5,628.6) -- (674.7,637.3) -- (677.8,649.0) -- (688.7,655.2) -- (691.7,661.3) -- (694.8,663.5) -- (699.0,669.6) -- (704.2,673.7) -- (704.9,678.1) -- (703.4,684.6) -- (697.8,694.4) -- (680.9,702.5) -- (680.6,708.7) -- (684.2,715.7) -- (680.4,725.9) -- (672.4,734.3) -- (672.7,739.6) -- (671.3,740.3) -- cycle;
\draw[mapstate] (671.3,740.3) -- (672.7,739.6) -- (672.4,734.3) -- (680.4,725.9) -- (684.2,715.7) -- (680.6,708.7) -- (680.9,702.5) -- (697.8,694.4) -- (704.4,681.8) -- (704.2,673.7) -- (699.0,669.6) -- (689.9,658.8) -- (790.3,652.5) -- (787.7,680.5) -- (779.3,681.2) -- (786.6,770.9) -- (784.6,774.2) -- (783.9,778.9) -- (788.2,788.4) -- (788.7,794.7) -- (785.1,804.6) -- (783.6,805.4) -- (780.0,812.3) -- (777.0,814.3) -- (778.0,818.2) -- (776.7,819.5) -- (776.1,826.2) -- (774.7,828.6) -- (776.6,831.7) -- (774.3,835.3) -- (775.4,841.5) -- (773.1,844.3) -- (764.3,848.7) -- (766.5,857.8) -- (765.8,859.1) -- (753.5,856.3) -- (748.3,856.8) -- (745.8,860.0) -- (745.5,864.2) -- (741.8,861.1) -- (741.1,863.7) -- (737.5,860.5) -- (735.0,856.5) -- (735.2,848.1) -- (732.9,840.3) -- (723.4,831.9) -- (712.6,825.5) -- (706.5,818.8) -- (707.5,811.1) -- (712.1,796.2) -- (709.5,794.4) -- (701.8,792.9) -- (699.0,794.5) -- (696.3,791.7) -- (693.8,784.5) -- (688.1,777.0) -- (674.1,764.0) -- (670.0,751.0) -- (671.3,740.3) -- cycle;
\draw[mapstate] (776.6,831.7) -- (774.7,828.6) -- (776.1,826.2) -- (776.7,819.5) -- (778.0,818.2) -- (777.0,814.3) -- (780.0,812.3) -- (783.6,805.4) -- (785.1,804.6) -- (788.9,792.0) -- (783.9,778.9) -- (784.6,774.2) -- (786.6,770.9) -- (779.3,681.2) -- (855.6,673.5) -- (866.5,773.6) -- (865.9,779.7) -- (868.4,781.9) -- (868.2,783.7) -- (858.0,789.8) -- (853.8,788.5) -- (851.0,789.1) -- (850.2,796.7) -- (847.8,798.7) -- (843.9,806.1) -- (841.5,806.9) -- (837.5,817.6) -- (835.8,818.9) -- (830.1,817.5) -- (825.0,812.6) -- (825.0,813.9) -- (822.4,815.1) -- (819.1,824.7) -- (810.8,821.7) -- (803.3,828.9) -- (796.0,825.0) -- (791.7,825.5) -- (789.5,824.4) -- (788.3,829.1) -- (786.0,826.8) -- (781.2,826.6) -- (779.6,831.3) -- (776.6,831.7) -- cycle;
\draw[mapstate] (358.4,862.6) -- (366.4,749.7) -- (463.1,755.1) -- (559.1,756.9) -- (565.8,761.8) -- (569.6,761.0) -- (571.2,762.2) -- (572.1,766.7) -- (570.7,766.8) -- (567.0,772.1) -- (567.5,775.5) -- (575.7,787.3) -- (580.3,789.9) -- (581.0,870.0) -- (469.6,868.5) -- (358.4,862.6) -- cycle;
\draw[mapstate] (745.5,864.2) -- (745.8,860.0) -- (748.3,856.8) -- (753.5,856.3) -- (765.8,859.1) -- (766.5,857.8) -- (764.3,848.7) -- (773.1,844.3) -- (775.4,841.5) -- (774.3,835.3) -- (777.1,831.3) -- (779.6,831.3) -- (780.6,826.7) -- (786.0,826.8) -- (788.3,829.1) -- (789.5,824.4) -- (791.7,825.5) -- (793.5,824.7) -- (803.3,828.9) -- (810.8,821.7) -- (819.1,824.7) -- (822.4,815.1) -- (825.0,813.9) -- (825.0,812.6) -- (830.1,817.5) -- (835.8,818.9) -- (837.5,817.6) -- (841.5,806.9) -- (843.9,806.1) -- (847.8,798.7) -- (850.2,796.7) -- (851.0,789.1) -- (853.8,788.5) -- (858.0,789.8) -- (868.2,783.7) -- (868.4,781.9) -- (865.9,779.7) -- (865.5,775.7) -- (867.9,772.2) -- (873.1,773.0) -- (877.8,772.3) -- (880.1,773.6) -- (884.8,780.0) -- (888.4,782.1) -- (897.0,783.6) -- (901.5,786.1) -- (905.6,783.7) -- (911.3,785.7) -- (915.7,784.7) -- (921.3,781.1) -- (924.8,780.7) -- (929.2,786.5) -- (934.8,790.4) -- (935.4,800.6) -- (943.7,811.7) -- (947.8,816.4) -- (953.6,820.0) -- (957.9,820.6) -- (948.0,831.3) -- (938.8,838.6) -- (933.6,848.2) -- (926.7,854.6) -- (910.1,863.5) -- (826.2,869.9) -- (786.1,874.3) -- (785.1,873.2) -- (778.2,873.1) -- (779.9,880.3) -- (737.8,883.3) -- (739.2,879.9) -- (745.0,878.1) -- (746.7,870.9) -- (745.5,864.2) -- cycle;
\draw[mapstate] (600.4,1020.6) -- (692.3,1017.7) -- (694.2,1020.6) -- (693.2,1027.6) -- (694.7,1028.6) -- (693.0,1032.5) -- (694.7,1037.6) -- (700.1,1043.4) -- (700.4,1045.4) -- (696.3,1051.2) -- (693.9,1057.4) -- (689.9,1060.9) -- (687.3,1065.2) -- (686.1,1070.3) -- (683.4,1075.0) -- (682.9,1081.7) -- (682.0,1083.9) -- (680.1,1084.4) -- (681.1,1090.7) -- (679.3,1094.9) -- (741.9,1091.7) -- (739.3,1102.2) -- (740.0,1106.7) -- (744.8,1112.1) -- (748.4,1121.0) -- (750.8,1122.1) -- (748.7,1123.2) -- (736.4,1120.0) -- (733.2,1117.1) -- (727.3,1116.3) -- (721.6,1125.8) -- (729.6,1129.6) -- (735.6,1128.1) -- (743.8,1123.3) -- (746.3,1125.2) -- (741.6,1129.7) -- (744.1,1132.5) -- (747.9,1133.3) -- (749.2,1132.8) -- (749.9,1129.4) -- (752.0,1127.1) -- (755.2,1127.4) -- (757.2,1135.9) -- (751.5,1138.1) -- (749.7,1140.0) -- (750.8,1141.0) -- (746.7,1142.1) -- (745.6,1144.3) -- (747.3,1147.3) -- (752.9,1151.5) -- (764.2,1154.0) -- (766.6,1157.2) -- (770.1,1158.7) -- (770.0,1161.0) -- (767.2,1162.9) -- (765.9,1166.1) -- (763.0,1163.8) -- (758.6,1167.9) -- (759.1,1164.5) -- (755.8,1160.0) -- (749.7,1156.2) -- (743.9,1155.4) -- (742.8,1152.3) -- (740.8,1150.8) -- (731.2,1148.3) -- (732.0,1150.5) -- (735.0,1152.1) -- (735.2,1155.7) -- (733.9,1161.8) -- (730.4,1164.9) -- (727.1,1159.3) -- (724.4,1157.9) -- (717.6,1159.2) -- (714.8,1164.0) -- (712.4,1164.9) -- (703.8,1162.9) -- (694.1,1159.7) -- (695.8,1158.0) -- (698.7,1158.4) -- (695.4,1153.2) -- (695.0,1149.1) -- (692.4,1151.2) -- (686.2,1149.8) -- (680.7,1142.7) -- (675.6,1142.8) -- (673.2,1139.6) -- (669.2,1141.1) -- (665.4,1145.1) -- (668.0,1148.1) -- (667.2,1149.1) -- (661.4,1150.7) -- (647.7,1149.6) -- (630.8,1143.1) -- (613.6,1143.8) -- (609.3,1145.5) -- (608.0,1144.4) -- (607.3,1142.3) -- (609.8,1140.7) -- (610.9,1138.3) -- (610.2,1135.9) -- (612.9,1130.7) -- (611.4,1121.9) -- (612.3,1113.7) -- (616.7,1103.5) -- (618.1,1095.0) -- (617.1,1089.9) -- (613.5,1084.6) -- (611.1,1077.3) -- (608.8,1075.3) -- (607.4,1066.8) -- (600.9,1059.0) -- (600.4,1020.6) -- cycle;
\draw[mapstate] (774.7,1135.2) -- (773.8,1136.5) -- (774.4,1130.6) -- (772.7,1125.8) -- (774.8,1130.4) -- (774.7,1135.2) -- cycle;
\draw[mapstate] (677.0,1152.3) -- (675.8,1152.9) -- (670.2,1149.8) -- (669.8,1148.4) -- (674.1,1147.0) -- (678.2,1149.7) -- (677.0,1152.3) -- cycle;
\draw[mapstate] (708.7,449.3) -- (714.4,448.0) -- (730.6,437.6) -- (736.4,435.6) -- (831.1,480.3) -- (834.8,490.0) -- (841.2,496.1) -- (845.1,493.4) -- (851.6,491.6) -- (854.6,503.3) -- (858.2,508.3) -- (867.5,505.9) -- (871.9,510.0) -- (869.5,517.1) -- (899.4,531.1) -- (920.0,596.1) -- (915.3,615.8) -- (913.9,633.1) -- (906.3,643.3) -- (901.0,647.3) -- (899.7,653.5) -- (901.4,660.6) -- (892.9,669.7) -- (855.9,675.7) -- (855.6,673.5) -- (787.7,680.5) -- (790.3,652.5) -- (784.2,621.4) -- (783.5,608.9) -- (785.0,592.2) -- (790.9,563.0) -- (795.0,555.2) -- (801.5,547.5) -- (787.2,540.2) -- (778.1,541.1) -- (775.2,547.5) -- (770.9,551.9) -- (770.2,556.4) -- (764.9,554.8) -- (762.7,550.4) -- (763.8,545.4) -- (759.1,546.2) -- (758.2,545.3) -- (759.6,537.2) -- (759.0,533.8) -- (755.8,531.3) -- (750.5,529.6) -- (750.8,527.2) -- (749.1,524.7) -- (740.0,522.8) -- (730.8,522.6) -- (723.1,518.9) -- (696.5,513.2) -- (693.2,507.0) -- (688.9,505.7) -- (688.3,504.4) -- (708.7,449.3) -- cycle;
\draw[mapstate] (512.1,415.4) -- (564.0,415.9) -- (564.2,401.5) -- (569.6,402.2) -- (571.7,403.9) -- (573.0,415.5) -- (576.3,424.2) -- (597.2,430.4) -- (598.4,432.6) -- (600.9,433.4) -- (607.0,432.0) -- (609.2,429.8) -- (618.8,429.8) -- (631.5,436.2) -- (633.9,442.2) -- (635.6,442.1) -- (636.8,440.1) -- (640.0,439.6) -- (644.3,440.8) -- (648.2,444.7) -- (657.1,449.7) -- (660.4,449.5) -- (668.9,444.0) -- (672.1,443.2) -- (674.1,443.5) -- (676.7,447.0) -- (693.3,445.6) -- (696.0,449.4) -- (698.4,450.0) -- (708.7,449.3) -- (698.6,477.0) -- (680.5,477.2) -- (660.1,492.8) -- (647.5,500.7) -- (643.3,499.6) -- (640.8,503.2) -- (639.1,503.6) -- (639.8,525.5) -- (629.4,532.5) -- (624.3,542.2) -- (624.5,544.6) -- (628.1,546.5) -- (630.2,550.6) -- (627.6,564.4) -- (628.6,577.5) -- (648.3,589.2) -- (654.4,596.8) -- (667.1,604.0) -- (668.9,605.9) -- (670.3,611.1) -- (671.7,622.3) -- (600.2,624.4) -- (528.9,624.6) -- (529.7,556.5) -- (527.7,553.7) -- (523.3,551.3) -- (519.3,543.9) -- (525.7,537.4) -- (527.3,532.0) -- (527.3,524.6) -- (521.7,505.1) -- (522.3,494.8) -- (520.9,491.0) -- (520.4,469.2) -- (514.0,446.8) -- (514.3,427.4) -- (512.1,415.4) -- cycle;
\draw[mapstate] (731.3,902.6) -- (710.7,904.1) -- (719.3,892.1) -- (719.2,889.3) -- (716.3,884.6) -- (648.7,887.5) -- (581.1,888.9) -- (580.3,789.9) -- (575.7,787.3) -- (567.5,775.5) -- (567.0,772.1) -- (570.7,766.8) -- (572.1,766.7) -- (571.2,762.2) -- (569.6,761.0) -- (565.8,761.8) -- (558.7,756.6) -- (555.8,749.5) -- (550.1,743.2) -- (547.2,734.2) -- (662.6,731.4) -- (671.3,740.3) -- (670.0,751.0) -- (674.1,764.0) -- (688.1,777.0) -- (693.8,784.5) -- (696.3,791.7) -- (699.0,794.5) -- (701.8,792.9) -- (709.5,794.4) -- (712.1,796.2) -- (707.5,811.1) -- (706.5,818.8) -- (712.6,825.5) -- (723.4,831.9) -- (732.9,840.3) -- (735.2,848.1) -- (735.0,856.5) -- (739.5,862.9) -- (741.4,863.7) -- (741.8,861.1) -- (745.5,864.2) -- (746.7,870.9) -- (745.0,878.1) -- (739.2,879.9) -- (737.2,883.7) -- (734.7,881.5) -- (735.3,886.2) -- (733.4,890.3) -- (734.3,892.1) -- (731.7,894.2) -- (732.9,898.6) -- (731.3,902.6) -- cycle;
\draw[mapstate] (780.9,937.0) -- (783.8,939.3) -- (780.0,1055.2) -- (787.2,1112.7) -- (777.7,1114.0) -- (770.5,1112.2) -- (760.3,1116.1) -- (757.0,1115.8) -- (753.3,1120.7) -- (750.8,1122.1) -- (748.4,1121.0) -- (744.8,1112.1) -- (740.0,1106.7) -- (739.3,1102.2) -- (741.9,1091.7) -- (679.3,1094.9) -- (681.1,1090.7) -- (680.1,1084.4) -- (682.0,1083.9) -- (682.9,1081.7) -- (683.4,1075.0) -- (686.1,1070.3) -- (687.3,1065.2) -- (689.9,1060.9) -- (693.9,1057.4) -- (696.3,1051.2) -- (700.1,1046.8) -- (700.1,1043.4) -- (696.2,1040.0) -- (693.0,1032.5) -- (694.6,1029.8) -- (693.1,1025.6) -- (694.2,1020.6) -- (692.6,1019.6) -- (692.2,1016.8) -- (693.3,1010.0) -- (691.4,1007.2) -- (691.9,1001.7) -- (689.3,996.3) -- (690.6,995.0) -- (689.4,992.8) -- (694.0,986.9) -- (693.5,980.8) -- (697.9,978.3) -- (697.0,975.3) -- (697.8,973.0) -- (700.6,970.5) -- (702.1,966.4) -- (706.0,964.1) -- (708.3,952.0) -- (714.8,945.2) -- (715.4,941.3) -- (780.9,937.0) -- cycle;
\draw[mapstate] (408.1,410.8) -- (512.1,415.4) -- (514.3,427.4) -- (514.0,446.8) -- (520.4,469.2) -- (520.9,491.0) -- (522.3,494.8) -- (521.7,505.1) -- (526.4,519.5) -- (527.3,532.0) -- (428.9,528.6) -- (330.8,521.2) -- (341.3,405.3) -- (408.1,410.8) -- cycle;
\draw[mapstate] (366.4,749.7) -- (369.2,712.1) -- (312.5,707.2) -- (319.6,632.0) -- (395.9,638.2) -- (473.7,642.2) -- (481.7,648.0) -- (486.6,650.0) -- (494.1,647.6) -- (506.8,648.4) -- (522.3,657.1) -- (525.5,661.4) -- (528.4,662.2) -- (530.3,665.5) -- (532.6,676.3) -- (538.4,688.1) -- (539.1,698.5) -- (542.3,701.8) -- (542.1,705.3) -- (543.7,707.2) -- (543.3,711.2) -- (544.3,712.1) -- (545.3,723.4) -- (544.5,728.8) -- (549.4,739.3) -- (550.1,743.2) -- (551.6,745.8) -- (554.1,747.0) -- (559.1,756.9) -- (463.1,755.1) -- (366.4,749.7) -- cycle;
\draw[mapstate] (901.4,660.6) -- (904.6,664.6) -- (914.9,669.2) -- (921.9,668.2) -- (951.2,643.5) -- (971.4,635.8) -- (981.7,698.3) -- (978.6,700.5) -- (978.2,702.1) -- (981.0,709.2) -- (977.5,736.5) -- (975.5,740.4) -- (968.0,748.6) -- (966.1,749.8) -- (962.8,749.2) -- (957.2,754.6) -- (954.1,762.2) -- (955.5,767.3) -- (952.4,769.8) -- (948.7,765.9) -- (946.9,768.0) -- (944.2,775.8) -- (945.1,781.8) -- (943.7,782.5) -- (940.6,789.0) -- (937.5,790.2) -- (929.2,786.5) -- (924.8,780.7) -- (921.3,781.1) -- (915.7,784.7) -- (911.3,785.7) -- (905.6,783.7) -- (901.5,786.1) -- (897.0,783.6) -- (888.4,782.1) -- (884.8,780.0) -- (880.1,773.6) -- (877.8,772.3) -- (871.1,772.9) -- (869.6,771.8) -- (866.5,773.6) -- (855.9,675.7) -- (892.9,669.7) -- (901.4,660.6) -- cycle;
\draw[mapstate] (581.1,888.9) -- (587.0,930.9) -- (586.2,996.6) -- (579.0,994.8) -- (563.1,985.0) -- (558.1,988.2) -- (553.4,987.7) -- (551.2,985.5) -- (545.2,988.5) -- (540.3,987.8) -- (534.0,990.1) -- (528.5,994.2) -- (519.8,988.9) -- (519.4,986.7) -- (515.0,988.6) -- (509.1,984.7) -- (505.0,988.3) -- (503.8,992.3) -- (502.3,992.8) -- (500.6,990.1) -- (500.9,987.0) -- (494.4,989.1) -- (489.1,985.7) -- (487.0,982.8) -- (480.5,987.0) -- (478.0,986.8) -- (477.2,983.3) -- (473.6,980.6) -- (472.8,976.6) -- (464.7,976.5) -- (461.6,978.8) -- (456.9,975.4) -- (453.0,975.9) -- (439.8,972.4) -- (437.1,966.1) -- (433.8,963.3) -- (432.3,965.7) -- (425.2,965.5) -- (418.2,959.3) -- (414.2,957.4) -- (417.8,885.2) -- (327.1,879.0) -- (328.7,860.2) -- (455.7,868.0) -- (581.0,870.0) -- (581.1,888.9) -- cycle;
\draw[mapstate] (528.9,624.6) -- (524.9,624.6) -- (526.5,630.7) -- (525.7,634.3) -- (527.7,635.7) -- (528.5,638.6) -- (523.7,652.6) -- (527.6,662.0) -- (523.4,659.7) -- (522.3,657.1) -- (518.7,654.6) -- (506.8,648.4) -- (494.1,647.6) -- (486.6,650.0) -- (481.7,648.0) -- (473.7,642.2) -- (395.9,638.2) -- (319.6,632.0) -- (330.8,521.2) -- (428.9,528.6) -- (527.3,532.0) -- (525.7,537.4) -- (519.3,543.9) -- (523.3,551.3) -- (527.7,553.7) -- (529.7,556.5) -- (528.9,624.6) -- cycle;
\draw[mapstate] (859.5,929.7) -- (780.3,936.1) -- (780.9,937.0) -- (715.4,941.3) -- (719.8,936.8) -- (721.2,933.3) -- (720.5,929.2) -- (725.4,918.4) -- (724.7,915.1) -- (726.0,912.4) -- (729.4,910.1) -- (729.7,906.9) -- (731.9,905.6) -- (731.2,903.3) -- (733.0,896.1) -- (731.7,894.2) -- (734.3,892.1) -- (733.4,890.3) -- (734.9,888.3) -- (734.6,883.4) -- (779.9,880.3) -- (778.2,873.1) -- (785.1,873.2) -- (786.1,874.3) -- (826.2,869.9) -- (910.1,863.5) -- (972.2,853.9) -- (971.4,863.8) -- (966.4,866.4) -- (963.1,873.9) -- (961.7,874.7) -- (959.0,873.3) -- (955.7,874.5) -- (948.7,881.7) -- (946.6,879.1) -- (944.0,880.2) -- (941.0,883.7) -- (938.7,883.9) -- (936.4,889.9) -- (932.6,891.4) -- (921.6,900.2) -- (911.0,904.2) -- (907.4,908.4) -- (906.2,913.5) -- (900.3,915.6) -- (899.5,925.6) -- (859.5,929.7) -- cycle;
\draw[mapstate] (586.2,996.6) -- (590.2,999.9) -- (595.4,998.8) -- (600.0,1000.1) -- (600.9,1059.0) -- (607.4,1066.8) -- (608.8,1075.3) -- (611.1,1077.3) -- (613.5,1084.6) -- (617.1,1089.9) -- (618.1,1095.0) -- (616.7,1103.5) -- (612.3,1113.7) -- (611.4,1121.9) -- (612.9,1130.7) -- (605.3,1142.2) -- (607.2,1146.9) -- (600.2,1147.8) -- (578.4,1158.9) -- (578.7,1157.6) -- (586.1,1152.6) -- (577.7,1152.7) -- (579.5,1148.5) -- (578.9,1144.9) -- (575.9,1144.8) -- (574.0,1147.7) -- (569.6,1146.8) -- (570.6,1153.3) -- (574.1,1159.4) -- (565.8,1167.2) -- (565.3,1170.6) -- (561.3,1175.0) -- (541.8,1187.4) -- (528.9,1193.1) -- (536.5,1187.6) -- (524.3,1186.6) -- (521.7,1188.9) -- (517.5,1184.2) -- (515.4,1184.5) -- (519.1,1192.9) -- (522.6,1194.2) -- (517.8,1197.6) -- (513.9,1198.5) -- (510.7,1195.4) -- (509.5,1203.1) -- (508.3,1204.1) -- (506.6,1205.5) -- (505.2,1202.9) -- (502.4,1205.2) -- (497.6,1205.8) -- (498.0,1209.0) -- (501.7,1207.8) -- (500.3,1211.9) -- (496.9,1215.9) -- (494.1,1216.8) -- (490.0,1216.2) -- (488.0,1217.4) -- (492.7,1223.9) -- (487.3,1236.9) -- (484.4,1237.4) -- (476.2,1231.7) -- (478.6,1238.3) -- (485.6,1240.4) -- (485.8,1245.0) -- (483.0,1250.7) -- (486.6,1269.4) -- (493.9,1285.7) -- (496.1,1285.8) -- (496.5,1287.2) -- (496.2,1289.8) -- (491.5,1290.5) -- (488.8,1293.2) -- (480.9,1288.6) -- (473.5,1286.2) -- (457.2,1283.0) -- (453.7,1280.1) -- (443.0,1276.1) -- (440.5,1273.6) -- (428.8,1269.0) -- (425.2,1256.7) -- (417.9,1245.0) -- (418.6,1240.6) -- (416.6,1234.0) -- (417.1,1226.0) -- (409.0,1218.6) -- (397.7,1201.0) -- (391.8,1194.7) -- (391.0,1188.0) -- (378.3,1161.2) -- (370.1,1152.3) -- (369.5,1150.0) -- (361.2,1142.8) -- (357.0,1137.0) -- (352.9,1135.5) -- (346.7,1136.1) -- (338.9,1135.0) -- (329.9,1131.5) -- (327.4,1131.5) -- (325.9,1133.6) -- (322.8,1134.8) -- (318.1,1135.1) -- (313.8,1138.9) -- (309.9,1146.2) -- (307.2,1154.7) -- (297.3,1162.3) -- (294.4,1161.9) -- (267.4,1144.6) -- (258.5,1136.5) -- (255.5,1132.2) -- (252.2,1125.1) -- (250.8,1119.7) -- (251.4,1114.3) -- (245.4,1096.4) -- (227.7,1079.3) -- (213.0,1061.6) -- (209.3,1059.1) -- (201.5,1046.1) -- (197.5,1043.6) -- (195.1,1036.3) -- (310.6,1048.5) -- (325.8,878.9) -- (417.8,885.2) -- (414.2,957.4) -- (418.2,959.3) -- (425.2,965.5) -- (432.3,965.7) -- (433.8,963.3) -- (437.1,966.1) -- (439.8,972.4) -- (453.0,975.9) -- (456.9,975.4) -- (461.6,978.8) -- (464.7,976.5) -- (472.8,976.6) -- (473.6,980.6) -- (477.2,983.3) -- (478.0,986.8) -- (480.5,987.0) -- (487.0,982.8) -- (489.1,985.7) -- (494.4,989.1) -- (500.9,987.0) -- (500.6,990.1) -- (502.3,992.8) -- (503.8,992.3) -- (505.0,988.3) -- (509.1,984.7) -- (515.0,988.6) -- (519.4,986.7) -- (519.8,988.9) -- (528.5,994.2) -- (534.0,990.1) -- (540.3,987.8) -- (545.2,988.5) -- (551.2,985.5) -- (553.4,987.7) -- (558.1,988.2) -- (563.1,985.0) -- (579.0,994.8) -- (586.2,996.6) -- cycle;
\draw[mapstate] (569.1,1168.1) -- (567.4,1168.4) -- (576.2,1160.6) -- (578.1,1160.6) -- (569.1,1168.1) -- cycle;
\draw[mapstate] (510.9,1205.7) -- (509.6,1205.9) -- (511.2,1203.8) -- (523.2,1197.1) -- (522.8,1198.8) -- (510.9,1205.7) -- cycle;
\draw[mapstate] (502.2,1215.2) -- (503.5,1210.9) -- (507.7,1206.4) -- (508.3,1208.1) -- (502.2,1215.2) -- cycle;
\draw[mapstate] (490.2,1238.1) -- (489.1,1240.3) -- (489.4,1237.0) -- (492.4,1229.5) -- (498.2,1219.8) -- (500.6,1218.2) -- (493.9,1228.9) -- (490.2,1238.1) -- cycle;
\draw[mapstate] (495.5,1282.2) -- (495.0,1284.0) -- (492.3,1275.5) -- (488.1,1256.5) -- (489.0,1242.0) -- (489.8,1257.3) -- (495.5,1282.2) -- cycle;
\draw[mapstate] (689.9,658.8) -- (688.7,655.2) -- (677.8,649.0) -- (676.3,645.5) -- (674.7,637.3) -- (675.9,630.6) -- (672.0,624.6) -- (668.9,605.9) -- (654.4,596.8) -- (648.3,589.2) -- (628.6,577.5) -- (627.6,564.4) -- (630.2,550.6) -- (628.1,546.5) -- (624.5,544.6) -- (624.3,542.2) -- (629.4,532.5) -- (639.8,525.5) -- (639.1,503.6) -- (640.8,503.2) -- (643.3,499.6) -- (647.5,500.7) -- (660.1,492.8) -- (680.5,477.2) -- (698.6,477.0) -- (688.3,504.4) -- (688.9,505.7) -- (693.2,507.0) -- (696.5,513.2) -- (723.1,518.9) -- (730.8,522.6) -- (740.0,522.8) -- (749.1,524.7) -- (750.8,527.2) -- (750.5,529.6) -- (755.8,531.3) -- (759.0,533.8) -- (759.6,537.2) -- (758.2,545.3) -- (759.1,546.2) -- (763.8,545.4) -- (762.7,550.4) -- (764.9,554.8) -- (770.2,556.4) -- (770.9,551.9) -- (775.2,547.5) -- (778.1,541.1) -- (787.2,540.2) -- (801.5,547.5) -- (795.0,555.2) -- (790.9,563.0) -- (785.0,592.2) -- (783.5,608.9) -- (784.2,621.4) -- (790.3,652.5) -- (689.9,658.8) -- cycle;

    #3
  \end{tikzpicture}}
\newcommand{\mapsz}{0.0196cm}

\usepackage{listings}
\lstset{language=Python,basicstyle=\ttfamily\footnotesize,breaklines=true,columns=fullflexible,keepspaces=true,frame=single,rulecolor=\color{gray!50},showstringspaces=false,upquote=true,xleftmargin=2pt,xrightmargin=2pt}
\usepackage{float}
